\section{Introduction}\label{sec:intro}

A time-projection readout turns a drift time into a coordinate. In a large TPC the drift coordinate is read along a track
many centimetres long and the local direction comes from many pad rows; in a micro-TPC ($\mu$TPC) the whole drift column is a
few millimetres to a few centimetres deep, each strip crossed by the track is given one drift time, and the inclination of the
segment is the slope of time against strip position~\cite{Alexopoulos2010,Bini2014}. The second form is attractive because it
recovers a good position for inclined tracks where the charge centroid degrades, and it is used in the ATLAS New Small Wheel
Micromegas~\cite{Iodice2015,Ntekas2016,Scholer2021}, in triple-GEM trackers~\cite{Alexeev2019,Farinelli2019,Flierl2018} and in
floating-strip Micromegas~\cite{Bortfeldt2014,Bortfeldt2015}.

Every one of these groups that compared the reconstructed inclination with the set angle of the chamber found the track too
steep, by a few degrees at \ang{10} shrinking to a fraction of a degree at \ang{40}~\cite{Bortfeldt2014,Iodice2015,Flierl2018,Ntekas2016}.
The published explanation is capacitive coupling: a strip at the edge of the cluster receives a copy of its neighbour's signal
at the neighbour's time, so the first and last strips carry the wrong depth~\cite{Bortfeldt2015,Iodice2015,Alexeev2019}. The
remedies are a cut on the charge ratio of the edge strips~\cite{Sessa2014,Rinnagel2019,Vanadia2015}, a shift of their assigned
position~\cite{Farinelli2019,Ntekas2016}, a subtraction of a fixed fraction of the neighbour pulse, sample by
sample~\cite{Loesel2017,Flierl2018,Herrmann2019}, or an angle-dependent calibration~\cite{Bortfeldt2014}. Two things are
missing from this record. No study has separated the coupling from the other processes that put charge of one depth on a
neighbouring strip, drift diffusion and the finite width of the induction footprint, and none has had a per-track external
reference for the angle: the comparison is always with the mounting angle of the chamber, through the peak or the mean of the
reconstructed distribution. Where the field's own simulations switch the coupling off they still find a bias, of \ang{0.2} to
\ang{4}~\cite{Ntekas2016} or \ang{7}~\cite{Scholer2021}, and leave it unexplained.

This paper argues that the question ``which process couples the strips'' is the wrong one for the angle. A single time per
channel, whatever its definition, is a functional of all the charge that reaches the channel, and the channels at the two ends of
a short column receive charge of the wrong depth from any lateral spread, by any mechanism. The bias is therefore a geometric
consequence of reading a short column one time per channel, its size is set by the lateral spread in units of the pitch and by the
number of channels in the column, and the mechanisms differ only in how they are to be modelled or removed. Pad TPCs and
liquid-argon TPCs have the same spread in units of the pitch but columns one or two orders of magnitude longer, and read a
charge-weighted time; directional and short-track detectors have the column but no reference that would reveal a few per cent
of slope. The strip $\mu$TPC is the one place where the effect is both large and measurable, and that is where it was found.

We test this account on a resistive-strip Micromegas with an unusually deep drift gap (\SI{28}{mm}), a fine pitch (\SI{0.78}{mm}),
sampled waveforms (\SI{60}{ns}) and a silicon telescope that gives the slope of every track independently of the chamber. The
detector, its calibration and its waveform reconstruction are described in two companion papers~\cite{paperI,paperII}; here it
serves as a test bed, with a forward model whose ingredients (footprint, diffusion, resistive spreading, electronics response)
can be switched on and off, and with 50\,000 track projections in two runs. \Cref{sec:formulation} sets out the formulation and
the scaling, \cref{sec:mechanisms} what the lateral spreads are and which observables separate them, \cref{sec:literature} what
each detector family has reported, \cref{sec:measurements} the measurements, \cref{sec:remedies} the remedies scored on the same
tracks, and \cref{sec:guidance} the consequences for the design and analysis of such readouts.

Terminology. We write $L$ for the drift length of the column, $p$ for the readout pitch, $v$ for the drift velocity, $\theta$
for the angle between the track and the drift direction, so that a track crosses $N = L\tan\theta/p$ strips, and $u$ for the
drift time of a charge from depth $z = v u$. ``Ratio'' means $\tan\theta_{\rm rec}/\tan\theta$ at fixed drift velocity, which is
also the factor by which the drift-time ladder is compressed; it is 1 for an unbiased reconstruction and larger than 1 when the
track reads too steep. ``Same-time'' sharing means charge appearing on a neighbour with the time structure of the strip's own
charge; ``delayed'' sharing means a later and longer signal, the signature of resistive spreading.
